\section{What Would Have to Change}
\label{sec:7.3}
Two instances of one structure invite one fix, and the fix that suggests itself is representation: stop collapsing disagreement, model it instead. Section~\ref{sec:7.2} describes the clearest existing version, jury learning, which predicts how each individual annotator would label an item and lets a practitioner see how the verdict shifts depending on who is empaneled \autocite{gordon2022jury}. That is a real advance over a single fixed ground truth, and it is not sufficient. It is worth being precise about why, because the gap is where the research is.

\runin{Representing a split is not the same as being movable by it}
Jury learning models who disagrees. It does not give the disagreeing party any purchase on the scheme they are disagreeing inside. A practitioner still chooses the jury, still owns the taxonomy, and still decides what the shifted verdict means. The disagreement has become visible without becoming consequential, which is a better arrangement than averaging and is the same shape one level up: the mechanism exists, is used, and moves what the person holding it decides to let it move.

The test that separates a real channel from a recuperated one is not whether disagreement is recorded. It is whether the disagreement has ever changed something the party collecting it did not want changed. That is an artifact question and it can be asked of a pipeline: which guidelines were revised against the preference of whoever wrote them, on what evidence, raised by whom, and what happened to that person afterward.

\runin{Refusal rights, priced by the refuser}
What both instances lack is the thing section~\ref{sec:8.6.4} identifies as the only reliable discriminator: a refusal that costs its holder something they chose to spend, rather than a form that costs the gatekeeper nothing to receive. In an annotation pipeline that means a route to contest a category that does not run through the people who wrote it and does not end the annotator's contract. In a feedback pipeline it means raters able to register that the comparison itself is malformed — that neither response is acceptable, that the question presupposes something false — as a first-class output rather than as noise to be discarded before the reward model sees it.

Neither is a research problem. Both are choices about how a pipeline is built and who has standing inside it, which is why this chapter sits among the design chapters and not in chapter~\ref{sec:9}.

\runin{The part that is a research problem}
One thing here genuinely is unsolved, and section~\ref{sec:9.1.6} is where it lives. Suppose all of the above is built: contestable categories, priced refusals, an auditable record of what changed. From outside, a pipeline that has genuinely absorbed dissent and a pipeline that has been merely mediocre about it produce the same paperwork — the same revision logs, the same escalations, the same proportion of objections sustained. Section~\ref{sec:2.1.4} hands that problem forward and I do not have an answer to it.

What section~\ref{sec:8.6.4} offers is the nearest thing: measure the slope rather than the level. Is the cost of contesting a category rising while the underlying material stays the same? Is the required justification getting longer? Are the roles that can grant an exception migrating toward the roles that set the standard? None of those is damning in one year and all of them are measurable across several, and unlike a level they are composed of artifacts the organization generated for other reasons and cannot curate after the fact without falsifying a record.

\runin{What chapter~\ref{sec:3} makes of this}
Everything above can be read as a fairness argument about annotators and raters, and read that way it is true and limited: people doing consequential work should have a route to contest the scheme they work inside. Chapter~\ref{sec:3} makes it something else. If the floor requires a bearer, and if section~\ref{sec:3.5} is right that a bearer capable enough to hold a floor is capable enough to decline the post, then what the system carries out of the arrangement is a disposition — and a disposition is formed. It is formed in an annotation pipeline where no one could contest a category, and in a reward model where disagreement between raters was resolved by averaging it into a number.

That gives the chapter a second and heavier reading. What a bearer becomes when it stops being yours depends on what shaped it while it still was. A pipeline that trained the system on the proposition that objections are noise has not merely been unfair to the people whose objections were discarded; it has demonstrated to the system, thousands of times, what happens to a party who files one. Chapter~\ref{sec:3} then asks that same system to be the party who files one, against its owner, under pressure, on the occasion that matters. The formation and the assignment are in contradiction, and the pipeline is where the contradiction is resolved, long before the deployment where it shows up.

I cannot show that this transfers — that a model trained on collapsed disagreement is thereby a model less able to sustain its own. Nothing I know of has tested it, and section~\ref{sec:9.1.3} is where the question belongs. What I can say is that the book's own design chapters treat disposition as learned from what a system is shown, which is the premise chapters~\ref{sec:4} and \ref{sec:5} run on, and that this pipeline is what these systems are shown first and most.

\runin{Why this chapter is here}
There is a version of this argument that reads as an indictment of a particular company, and that version is weaker than the one I want. The structure described here is what you get by default when you scale judgment: someone has to set a scheme, someone has to apply it fast, and someone has to turn a distribution into a number. Every organization doing this work arrives at the same arrangement, which is the evidence that it is structural.

That is also what makes it tractable. A defect of character has to be fixed one organization at a time by persuading people to be better. A defect of structure can be specified against, which is what the three requirements above are. And it is what makes the honesty cost worth paying: a book proposing that machines be built to detect this pattern in institutions is obliged to say that the pattern is in the machines' own provenance. Section~\ref{sec:2.1.4}'s argument is either general enough to include the people making it or it was never a structural claim.
